\section{Three Boundary-Primitive Cases}
\label{sec:evidence}

The three cases in this section arose in different subsystems and have
different maturity. Their common feature is not that they record data. Each
names a contact surface where participants previously relied on implicit human
coordination.

\begin{table}[ht]
\centering
\small
\begin{tabular}{p{0.16\linewidth}p{0.30\linewidth}p{0.38\linewidth}}
\toprule
Candidate & Contact surface & Implicit mechanism replaced \\
\midrule
Xinfa Atlas & project reality $\leftrightarrow$ human/agent understanding &
Repeated repository interpretation, private prompts, hand-written summaries,
and divergent human/agent views \\
Kungfu Episode & agent action $\leftrightarrow$ real-world result &
Sessions, logs, process exit, and human reconstruction of which facts,
artifacts, receipts, and proof belong to one unit of work \\
Release Passport & reviewed source/build $\leftrightarrow$ released artifact &
Maintainer memory about what was promoted, which checks ran, which artifacts
exist, and what users can verify \\
\bottomrule
\end{tabular}
\caption{Three candidate primitives and the implicit boundaries they mediate.}
\label{tab:cases}
\end{table}

\subsection{Xinfa Atlas: reality to shared understanding}

A Xinfa Atlas is defined as the smallest portable, reproducible, and
verifiable semantic closure compiled at a declared source and policy cut, from
which equivalent human and agent views can be derived without reinterpreting
the unbounded source corpus \cite{xinfaAtlas2026}. It binds scope, provenance,
verification state, gaps, conflicts, routes, and deterministic identity while
remaining a representation rather than source authority.

The boundary pressure is cognitive. A repository may be readable to a
maintainer because that person carries years of implicit context. New humans,
agents, documentation sites, IDEs, and task runners otherwise reconstruct
different project models. A context window or archive is too narrow as the
primitive: it names a transport or one reader's selection. Atlas names the
cut-bound closure from which bounded Human Views, Agent Task Charts, and GUI
views can be derived under the same root \cite{xinfaProjections2026}.

The deletion test is strong but not conclusive. Without Atlas, the system
either repeats the full closure definition, elevates a storage pack into the
user object, or allows each projection to become a separate fact source. The
narrower alternative is retrieval over raw repositories. That alternative
would falsify Atlas if it delivered equivalent provenance, portability,
bounded cost, gap visibility, deterministic identity, and cross-reader parity
without recreating Atlas semantics.

As of the evidence cut used by this paper, Atlas identity and bounded
projections are accepted and implemented in Kungfu's Xinfa subsystem. Adoption
and measurable context-cost reduction remain open evaluation questions.

\subsection{Kungfu Episode: action to consequence}

A Kungfu Episode is a bounded causal unit of actual work whose facts,
artifacts, dependencies, receipts, and verification roots can be inspected,
sealed, exported, recovered, and used to support decisions
\cite{kungfuEpisode2026}. The object is deliberately not a chat session,
process, or bag of logs. Those coordinates end for technical reasons, while
real responsibility may span runs and preserve useful partial work after a
failure.

The boundary pressure is causal and accountable. A successful tool call does
not prove visibility, durability, projection, or completion. A process exit
does not state which artifacts belong to the work. A trace can show calls
without defining the authoritative unit that export, fsck, recovery, and later
trust assessment should address. Episode collects those relations under one
bounded lifecycle while keeping Claim distinct from Proof and Decision.

The deletion test asks whether runtime consumers can still export, verify,
recover, and reason about one unit of actual work without reconstructing it
from sessions, logs, traces, and external files. A narrower event-group label
would falsify the need for Episode if it supported the same causal closure,
identity, lifecycle, receipts, and recovery semantics.

Episode is the central product vocabulary and an accepted architectural
object, but its implementation evidence is mixed. Core causal-segment work is
active while some sealed identity, cross-provider equivalence, and bundle
contracts remain proposed or partial. This paper therefore treats Episode as
a strong, actively dogfooded candidate rather than a completed universal
standard.

\subsection{Buildchain Release Passport: build activity to released truth}

Buildchain Release Passport is a durable release responsibility record for
artifacts that users or agents depend on \cite{buildchainPassport2026}. Its
first-read artifact, \texttt{buildchain.release.json}, binds reviewed source,
release line and channel, exact tags, artifact hashes, checks, surface impact,
package publication, residual evidence, and transaction state.

The boundary pressure comes after CI says green. Build systems can compile,
test, and upload while leaving the recipient unable to answer which reviewed
state was released, what exact artifact was produced, which floating channel
now points to it, or how to verify and recover the transaction. Historically a
maintainer reconstructs this story from workflow pages, tags, package
registries, and memory.

The deletion test is whether an independent human or agent can reconstruct the
same release responsibility record from ordinary provider state at comparable
cost and without private maintainer context. The narrower alternative is a
provenance attestation alone. That alternative is sufficient only if it also
captures promotion, compatibility impact, package-set and channel state,
residual risk, and recovery responsibility rather than artifact origin alone.

This is the most mature case in the paper. Buildchain stable release
\texttt{v2.12.7} publishes the passport and its supporting evidence as GitHub
Release assets. Maturity in one ecosystem does not prove generality, but it
does show that the object can bear real release operations rather than remain
only a diagram.

\subsection{Shared structure and bounded claim}

The cases support a common pattern:

\begin{equation}
\text{implicit coordination}
\xrightarrow{\Delta P}
\text{recurring boundary failure}
\xrightarrow{X}
\text{inspectable mediation}.
\end{equation}

They do not yet establish that the pattern is necessary for all primitives or
that any case will reach historical significance. The evidence supports a
candidate-generation method and three non-trivial implementation cases. Time,
independent adoption, and adversarial deletion tests must carry the stronger
claim.
